\GalSourcePageBoundary{083}{L0082}{54}{54}
ainsi:
\[
\begin{array}{*{6}{l}}
  a_{0,0}, & a_{0,1}, & a_{0,2}, & a_{0,3}, & \dots, & a_{0,p-1},\\
  a_{1,0}, & a_{1,1}, & a_{1,2}, & a_{1,3}, & \dots, & a_{1,p-1},\\
  a_{2,0}, & a_{2,1}, & a_{2,2}, & a_{2,3}, & \dots, & a_{2,p-1},\\
  \dots,  & \dots,  & \dots,  & \dots,  & \dots, & \dotfill, \\
  a_{p-1,0}, & a_{p-1,1}, & a_{p-1,2}, & a_{p-1,3}, & \dots, & \rlap{$a_{p-1,p-1},$}
\end{array}
\]
chaque ligne horizontale et chaque ligne verticale étant un système
de lettres conjointes.

Si l'on permute entre elles les lignes horizontales, le groupe
que l'on obtiendra, étant primitif et de degré premier, ne devra
contenir que des substitutions de la forme
\[
\GalPrintedError{(a_{k_{1},k_{2}},\ a_{mk_{1}+n k_{2}})}{W10-PE003},
\]
les indices étant pris relativement au module~$p$.

Il en sera de même pour les lignes qui ne pourront donner que
des substitutions de la forme
\[
(a_{k_{1},k_{2}},\ a_{k_{1}, qk_{2}+r}).
\]
Donc enfin toutes les substitutions du groupe~$H$ seront de la
forme
\[
\GalPrintedError{(a_{k_{2},k_{3}},\ a_{m_{1} k_{1}+n_{1}, m_{2} k_{2}+n_{2}})}{W10-PE004}.
\]
Si un groupe~$G$ se partage en $n$~groupes conjugués à celui que
nous venons de décrire, toutes les substitutions du groupe~$G$ devront
transformer les unes dans les autres les substitutions circulaires
du groupe~$H$, qui sont toutes écrites comme il suit:
\[
\tag{a}
(a_{k_{1},k_{2}},\ \dots,\ a_{k_{1}+\alpha_{1}, k_{2}+\alpha_{2}},\ \dots).
\]
Supposons donc que l'une des substitutions du groupe~$G$ se forme
en remplaçant respectivement
\begin{align*}
k_{1} & \quad\text{par}\quad \phi_{1} (k_{1}, k_{2}),\\
k_{2} & \quad\text{par}\quad \phi_{2} (k_{1}, k_{2}).
\end{align*}
Si, dans les fonctions $\phi_{1}$,~$\phi_{2}$, on substitue pour $k_{1}$~et~$k_{2}$ les valeurs
$k_{1} + \alpha_{1}$, $k_{2} + \alpha_{2}$, il devra venir des résultats de la forme
\[
\phi_{1} + \beta_{1},\quad \phi_{2} + \beta_{2},
\]
